\scanpage{240}
\mspage{114}
% archive sheet number: 239

On a $i_X\in C$ par hypothèse, donc $j\in C$ (stabilité de $C$ par coch. de base), on a $\lambda:X\longrightarrow Y\amalg X$ dans $C$ car $X$ dans $\unsure{\mathcal{M}|M}_C$, donc le composé $i_f=j\circ\lambda$ est dans $C$ (stabilité de $C$ par composition) --- on gagne !

% slanted
\margin{\uline{Commentaires}}

Pour l'équivalence de $a)$ $b)$ $c)$ dans le th. 8.4, $1^\circ$, on n'a \deleted{\unsure{plus}} pas eu à faire appel aux \deleted{conditions} développements de PS, qui servent \deleted{à \unsure{établir}} seulement à établir l'équivalence des trois conditions en question avec la condition ``\unsure{modélisante|anodinisante}'' $d)$. Pour résumer :

$1^\circ)$ les implications

\[
a\Rightarrow(b\Leftrightarrow c\Leftrightarrow d)
\]

sont valables sous les seules conditions

\[
W(1,2\mathrm{bis},3)
\]

et \deleted{\unsure{l'équi--}} pour l'équivalence de $a$ aux autres conditions, il suffit de

\[
W(1\mathrm{bis},2\mathrm{bis},3)
\]

(où $W(1,\mathrm{bis})$ renforce $W(1)$ par la condition des fact. directs, tandis que $W(2\mathrm{bis})$ affaiblit $W(2)$, en exigeant seulement que les morphismes $W$-asph. soient dans $\unsure{W^{\circ}|W}$)

\deleted{\unsure{$2^\circ)$}} Les conclusions de 8.4. $3^\circ)$, relatives au cas où $A$ est $W$-asphérique, sont valables également sous les conditions $W(1,2\mathrm{bis},3)$
\inserted{\deleted{\unsure{sous|sans}} plus}
